\chapter{Chapter09 — Level Weight Asymptotics}
%%% generated by the blueprint pipeline from FABL.Chapter09.LevelWeightAsymptotics
\section{Overview}

Declarations of the Lean module \texttt{FABL.Chapter09.LevelWeightAsymptotics}.

\section{Declarations}

\begin{theorem}[A bound for the logarithm of four]
\lean{FABL.log_four_lt_seven_fifths}
\label{FABL.log_four_lt_seven_fifths}
\leanok
\difficulty{3}
The natural logarithm of $4$ is less than $7/5$.
\end{theorem}

\begin{theorem}[Numerical exponent bound]
\lean{FABL.exercise9_19_exponent_bound}
\label{FABL.exercise9_19_exponent_bound}
\leanok
\uses{FABL.log_four_lt_seven_fifths}
\difficulty{2}
The real numbers satisfy $\frac{21}{20} \le \frac{8}{5}-\frac{373}{1000}\log 4$.
\end{theorem}

\begin{theorem}[Low-degree Fourier weight bound]
\lean{FABL.exercise9_19_quantitative}
\label{FABL.exercise9_19_quantitative}
\leanok
\uses{FABL.exercise9_19_exponent_bound, FABL.fourierWeightAtMost, FABL.fourierWeightAtMost_le_invPow_mul_rpow, FABL.mean}
\difficulty{6}
Let $n,k$ be nonnegative integers, let $f:\{-1,1\}^n\to\mathbb R$ take only the values
$0$ and $1$, and suppose its uniform mean is $\alpha$, where $0<\alpha\le 1$. If $k\le
\frac{373}{1000}\log(1/\alpha)$, then the sum of the squares of the Fourier coefficients
of $f$ over subsets of at most $k$ coordinates is at most $\alpha^{21/20}$.
\end{theorem}

\begin{theorem}[A power is little-o of its argument]
\lean{FABL.rpow_twentyone_twentieth_isLittleO}
\label{FABL.rpow_twentyone_twentieth_isLittleO}
\leanok
\difficulty{5}
As $\alpha\to0$ through positive real values, the function $\alpha\mapsto\alpha^{21/20}$
is little-$o$ of the function $\alpha\mapsto\alpha$; equivalently, their ratio tends to
zero.
\end{theorem}

\begin{theorem}[Uniform asymptotic Fourier weight bound]
\lean{FABL.exercise9_19_uniform_littleO}
\proofsource{Analysis-of-Boolean-Functions-by-Ryan-ODonnell}{pdf-ffbe9c7b3306-p273-b006}
\label{FABL.exercise9_19_uniform_littleO}
\leanok
\uses{FABL.exercise9_19_quantitative, FABL.fourierWeightAtMost, FABL.mean, FABL.rpow_twentyone_twentieth_isLittleO}
\difficulty{4}
For every $\varepsilon>0$, there is a right-hand neighborhood of $0$ such that, for
every $\alpha$ in that neighborhood with $\alpha\le1$, every natural number $n$, every
function $f:\{-1,1\}^n\to\mathbb R$ taking only the values $0$ and $1$ whose uniform
mean is $\alpha$, and every natural number $k\le(373/1000)\log(1/\alpha)$, the sum of
the squared Fourier coefficients of $f$ over subsets of at most $k$ coordinates is at
most $\varepsilon\alpha$.
\end{theorem}
